\begin{abstract}
\noindent Este artigo discute se a dependência de transferências, indicador usado para avaliar a autonomia fiscal dos municípios brasileiros, mede o que promete. A questão nasce de uma distorção contábil: a cota-parte do ICMS é registrada como transferência do estado, embora seja repartida conforme a riqueza produzida em cada município, e por isso recebe o mesmo nome de repasses que vêm de fora, como o Fundo de Participação dos Municípios (FPM). Para examinar essa distorção, o estudo compara Barreiras e Camaçari, na Bahia, que recebem o mesmo FPM, mas têm economias opostas, e contrasta esses casos com os governos locais da Hungria, um Estado unitário que recentralizou suas finanças. A pesquisa atualiza até 2025 e corrige a base de um trabalho anterior e combina análise descritiva com regressões que testam o efeito \textit{flypaper}, segundo o qual o gasto local reage mais às transferências do que à renda dos moradores. Os resultados mostram que o indicador convencional ordena corretamente os dois municípios, mas superestima a dependência de Camaçari: sem a cota do ICMS gerada por sua indústria, a parcela de sua receita que vem de fora cai, em média, de 61\% para 27\%. Na Hungria, a dependência caiu enquanto os governos locais perdiam competências. Já o teste do \textit{flypaper} não consegue separar a resposta do gasto às transferências da simples regra de gastar o que se arrecada. Ao expor os limites dessa medida, o artigo contribui para avaliar melhor a autonomia municipal, às vésperas da substituição do ICMS pela reforma tributária, e oferece uma base de dados aberta e reprodutível.

\vspace{0.5em}
\noindent\textbf{Palavras-chave:} transferências intergovernamentais; efeito \textit{flypaper}; ilusão fiscal; autonomia fiscal; cota-parte do ICMS; Hungria.\\
\noindent\textbf{Classificação JEL:} H71; H72; H77.\\
\noindent\textbf{Dados e códigos:} \url{https://github.com/joaogabrielmarecogalhas9154-blip/P2-FSP-232045821}
\end{abstract}
